\documentclass[../../../topic_linear-algebra]{subfiles}

\usepackage{xr-hyper}
\externaldocument{../../../.tex_intermediates/topic_linear-algebra}

\begin{document}

\sectionline
\section{部分空間の和の次元}
\marginnote{\refbookR p54〜55 \\ \refbookA p103 \\ \refbookC p39〜41 \\ \refbookF 235〜236}

部分空間の和（和空間）を考える上で、部分空間の共通部分が重要だと述べたのは、和空間の次元を考える際に、共通部分の次元が関わってくるからである。

\br

次のように、$V$の部分空間$U, W$の和空間の次元（基底をなすベクトルの個数）を考えるときは、共通部分の基底が重複して数えられてしまうため、共通部分の次元を引く必要がある。

% \refbookR 定理5.3
\begin{theorem}{部分空間の和の次元}{dimension-of-sum-of-subspaces}
  $V$を$K$上の有限次元線形空間とし、$U, W$を$V$の部分空間とする。
  
  このとき、次が成り立つ。
  \begin{equation*}
    \dim(U + W) = \dim U + \dim W - \dim(U \cap W)
  \end{equation*}
\end{theorem}

\begin{proof}
  $\dim U = n, \, \dim W = m$とする。

  \br

  $U \cap W$の基底$\mathcal{V} = \{ \vb*{v}_1,\dots, \vb*{v}_d \}$をとる。

  これを\thmref{thm:basis-extension}に基づいて、$U$の基底
  \begin{equation*}
    \mathcal{V} \cup \{ \vb*{u}_1, \dots, \vb*{u}_{n-d} \}
  \end{equation*}
  に延長する。

  同様に、$\mathcal{V}$を$W$の基底
  \begin{equation*}
    \mathcal{V} \cup \{ \vb*{w}_1, \dots, \vb*{w}_{m-d} \}
  \end{equation*}
  に延長する。

  \br

  このとき、$\vb*{v}_1, \dots, \vb*{v}_d, \vb*{u}_1, \dots, \vb*{u}_{n-d}, \vb*{w}_1, \dots, \vb*{w}_{m-d}$が$W_1 + W_2$の基底になることを示す。

  \begin{subpattern}{\bfseries $U+W$を生成すること}
    $\vb*{u} \in U, \, \vb*{w} \in W$とすると、それぞれ基底の線形結合で表すことができる。
    \begin{align*}
      \vb*{u} & = \sum_{i=1}^d a_i \vb*{v}_i + \sum_{j=1}^{n-d} b_j \vb*{u}_j \\
      \vb*{w} & = \sum_{i=1}^d c_i \vb*{v}_i + \sum_{k=1}^{m-d} d_k \vb*{w}_k
    \end{align*}

    \br

    $U+W$の任意の元は、$\vb*{u} + \vb*{w}$と書けるので、
    \begin{equation*}
      \vb*{u} + \vb*{w} = \sum_{i=1}^d (a_i + c_i) \vb*{v}_i + \sum_{j=1}^{n-d} b_j \vb*{u}_j + \sum_{k=1}^{m-d} d_k \vb*{w}_k
    \end{equation*}
    となり、$\{ \vb*{v}_1, \dots, \vb*{v}_d, \vb*{u}_1, \dots, \vb*{u}_{n-d}, \vb*{w}_1, \dots, \vb*{w}_{m-d} \}$の線形結合で表せる。 $\qed$
  \end{subpattern}

  \begin{subpattern}{\bfseries 線型独立であること}
    $\vb*{v}_1, \dots, \vb*{v}_d, \vb*{u}_1, \dots, \vb*{u}_{n-d}, \vb*{w}_1, \dots, \vb*{w}_{m-d}$が線型独立であることを示すために、次のような線形関係式を考える。
    \begin{equation*}
      \sum_{i=1}^d c_i \vb*{v}_i + \sum_{j=1}^{n-d} c_{d+j} \vb*{u}_j + \sum_{k=1}^{m-d} c_{d+n-d+k} \vb*{w}_k = \vb*{0}
    \end{equation*}

    ここで、$c_i \in K$はスカラーである。

    \br

    この式を$U$と$W$の基底の線型結合として考えると、$U$の基底$\vb*{v}_i, \, \vb*{u}_j$に関する部分と$W$の基底$\vb*{v}_i, \, \vb*{w}_k$に関する部分がそれぞれ線形独立であるため、結局どの項においても$c_i = 0$である必要がある。

    \br

    よって、$\vb*{v}_1, \dots, \vb*{v}_d, \vb*{u}_1, \dots, \vb*{u}_{n-d}, \vb*{w}_1, \dots, \vb*{w}_{m-d}$は線型独立である。 $\qed$
  \end{subpattern}

  \br

  以上より、$\vb*{v}_1, \dots, \vb*{v}_d, \vb*{u}_1, \dots, \vb*{u}_{n-d}, \vb*{w}_1, \dots, \vb*{w}_{m-d}$は$U + W$の基底をなす。

  \br

  この基底をなすベクトルの個数（次元）について考えると、
  \begin{align*}
    \dim(U + W) & = d + (n - d) + (m - d) \\
                & = n + m - d
  \end{align*}
  ここで、$d = \dim(U \cap W)$なので、
  \begin{equation*}
    \dim(U + W) = \dim U + \dim W - \dim(U \cap W)
  \end{equation*}
  と書き換えられ、目的の式が得られた。 $\qed$
\end{proof}

\subsection{複数の部分空間の場合}

より一般に、複数の部分空間$W_1, \ldots, W_n$の和の次元も考えることができる。

$n=2$の場合が、先ほど示した
\begin{equation*}
  \dim(W_1 + W_2) = \dim W_1 + \dim W_2 - \dim(W_1 \cap W_2)
\end{equation*}
であり、数学的帰納法を用いることにより、次の定理として一般化できる。

% \refbookR 定理5.4
\begin{theorem*}{部分空間の和の次元（一般の場合）}
  $V$を$K$上の有限次元線形空間とし、$W_1, \ldots, W_n$を$V$の部分空間とする。

  このとき、次が成り立つ。
  \begin{multline*}
    \dim(W_1 + \cdots + W_n) \\
    = \sum_{i=1}^n \dim W_i - \sum_{i=2}^n \dim \left( (W_1 + \cdots + W_{i-1}) \cap W_i \right)
  \end{multline*}
\end{theorem*}

\begin{proof}
  $n = 2$のとき、\thmref{thm:dimension-of-sum-of-subspaces}より、定理は成り立つ。
  
  \br
  
  $n = k \geq 2$のとき、定理が成り立つと仮定する。

  このとき、$U_j = W_1 + \cdots + W_j$とおくと、$U_j$は$V$の部分空間であるから、\thmref{thm:dimension-of-sum-of-subspaces}より、
  \begin{align*}
    \dim(W_1+\cdots+W_{k+1})
    &= \dim(U_k + W_{k+1}) \\
    &= \dim U_k + \dim W_{k+1} - \dim(U_k \cap W_{k+1})
  \end{align*}
  
  ここで、帰納法の仮定より、
  \begin{equation*}
    \dim U_k
=\sum_{i=1}^k \dim W_i
 -\sum_{i=2}^k
   \dim\bigl((W_1+\cdots+W_{i-1})\cap W_i\bigr)
  \end{equation*}
  であるから、これを代入すると、
  \begin{align*}
    &\quad \dim(W_1+\cdots+W_{k+1}) \\
    &= \sum_{i=1}^k \dim W_i -\sum_{i=2}^k \dim\bigl((W_1+\cdots+W_{i-1})\cap W_i\bigr) \\
    &\qquad + \dim W_{k+1} - \dim(U_k \cap W_{k+1}) \\
    &= \sum_{i=1}^k \dim W_i -\sum_{i=2}^k \dim\bigl((W_1+\cdots+W_{i-1})\cap W_i\bigr) \\
    &\qquad + \dim W_{k+1} - \dim(\bigl(W_1 + \cdots + W_k\bigr) \cap W_{k+1})
  \end{align*}
  
  ここで現れた和の項は、次のようにまとめることができる。
  \begin{equation*}
    \sum_{i=1}^{k+1} \dim W_i = \sum_{i=1}^{k} \dim W_i + \dim W_{k+1}
  \end{equation*}
  
  また、共通部分の項は、次のようにまとめることができる。
  \begin{multline*}
    \sum_{i=2}^{k+1} \dim\bigl((W_1+\cdots+W_{i-1})\cap W_i\bigr) \\
= \underbrace{\sum_{i=2}^{k} \dim\bigl((W_1+\cdots+W_{i-1})\cap W_i\bigr)}_{\text{もともとの和}} \\
+ \underbrace{\dim\bigl((W_1+\cdots+W_k)\cap W_{k+1}\bigr)}_{i=k+1\text{項}}
  \end{multline*}
  
  よって、
  \begin{multline*}
    \dim(W_1+\cdots+W_{k+1}) \\
    = \sum_{i=1}^{k+1} \dim W_i - \sum_{i=2}^{k+1} \dim\bigl((W_1+\cdots+W_{i-1})\cap W_i\bigr)
  \end{multline*}
  となり、$n = k + 1$のときも定理が成り立つ。
  
  数学的帰納法により、定理は任意の$n \geq 2$について成り立つことが示された。 $\qed$
\end{proof}

\end{document}
